% ==============================================================================
% Chapter 12: Spatial Transcriptomics
% ==============================================================================
\chapter{Spatial Transcriptomics}
\label{ch:spatial}

\begin{sourcebox}
The field is anchored here by the original array-based spatial transcriptomics report \cite{stahl2016}. Resolution, area, plex, sensitivity, and analysis-software claims for commercial systems require current product documentation and must not be compared as if they measured identical quantities.
\end{sourcebox}

\begin{keyconcept}
Spatial transcriptomics measures gene expression while preserving the physical
location of each measurement within a tissue section. By retaining spatial
context---lost in both bulk and dissociated single-cell RNA-seq---these
technologies reveal how cells are organized, how they communicate with their
neighbors, and how tissue architecture relates to function and disease.
\end{keyconcept}

% ---------------------------------------------------------------------------
\section{Why Spatial Context Matters}
\label{sec:spatial:why}

The human body is not a bag of cells: it is a precisely organized structure
where cellular function depends critically on spatial position. Conventional
scRNA-seq requires tissue dissociation, which destroys the spatial
relationships between cells. This loss has several consequences:

\begin{itemize}[nosep]
  \item \textbf{Tissue architecture is invisible:} The layered organization of
    the cortex, the crypt--villus axis of the intestine, and the zonal
    structure of the liver cannot be observed in dissociated data.
  \item \textbf{Cell neighborhoods are lost:} A tumor cell adjacent to a
    cytotoxic T cell behaves differently from one surrounded by
    immunosuppressive macrophages. Spatial context determines these
    interactions.
  \item \textbf{Gradients cannot be detected:} Morphogen gradients during
    development, oxygen gradients in tumors, and signaling gradients in wound
    healing require spatial measurement.
  \item \textbf{Rare niches are obscured:} Stem cell niches, germinal centers,
    and tertiary lymphoid structures are defined by their spatial organization
    as much as by their cell type composition.
  \item \textbf{Dissociation introduces artifacts:} Enzymatic digestion
    activates stress-response genes and selectively loses fragile cell types
    (neurons, adipocytes, syncytiotrophoblasts).
\end{itemize}

Spatial transcriptomics solves these problems by measuring gene expression
\emph{in situ}---within intact tissue sections.


% ---------------------------------------------------------------------------
\section{Categories of Spatial Technologies}
\label{sec:spatial:categories}

Spatial transcriptomics technologies can be broadly divided into three
categories based on their detection mechanism: sequencing-based methods that
capture RNA at defined spatial positions, imaging-based methods that visualize
individual RNA molecules, and in situ sequencing methods that read out RNA
sequences directly in tissue.

% TikZ diagram: Comparison of spatial transcriptomics approaches
\begin{figure}[htbp]
\centering
\begin{tikzpicture}[
  node distance=0.5cm and 1.0cm,
  every node/.style={font=\small},
  category/.style={draw, rounded corners, fill=MidnightBlue!15,
    minimum width=4.0cm, minimum height=0.9cm, align=center,
    font=\small\bfseries},
  platform/.style={draw, rounded corners, fill=white,
    minimum width=3.8cm, minimum height=0.7cm, align=center,
    font=\scriptsize},
  arrow/.style={-{Stealth[length=5pt]}, thick, MidnightBlue!60},
  label/.style={font=\scriptsize\itshape, text=gray!60!black},
]

% Categories
\node[category] (seq) {Sequencing-based\\(capture spots/beads)};
\node[category, right=2.5cm of seq] (img) {Imaging-based\\(in situ hybridization)};
\node[category, right=2.5cm of img] (iss) {In situ sequencing};

% Sequencing-based platforms
\node[platform, below=0.6cm of seq] (visium) {10x Visium (55\,$\mu$m spots)};
\node[platform, below=0.3cm of visium] (visiumhd) {10x Visium HD (2\,$\mu$m)};
\node[platform, below=0.3cm of visiumhd] (slideseq) {Slide-seq V2 (10\,$\mu$m beads)};
\node[platform, below=0.3cm of slideseq] (stereo) {Stereo-seq (nanoscale)};

% Imaging-based platforms
\node[platform, below=0.6cm of img] (merfish) {MERFISH / Vizgen (100s--1000s genes)};
\node[platform, below=0.3cm of merfish] (xenium) {10x Xenium (100s genes)};
\node[platform, below=0.3cm of xenium] (seqfish) {seqFISH+ (10,000+ genes)};
\node[platform, below=0.3cm of seqfish] (cosmx) {CosMx / NanoString};

% In situ sequencing
\node[platform, below=0.6cm of iss] (issplat) {ISS / BaristaSeq};
\node[platform, below=0.3cm of issplat] (starmap) {STARmap PLUS};
\node[platform, below=0.3cm of starmap] (exseq) {ExSeq (expansion seq.)};

% Arrows from categories to first platform
\draw[arrow] (seq.south) -- (visium.north);
\draw[arrow] (img.south) -- (merfish.north);
\draw[arrow] (iss.south) -- (issplat.north);

% Annotations at bottom
\node[label, below=0.5cm of stereo] {Genome-wide; lower resolution};
\node[label, below=0.5cm of cosmx] {Targeted; subcellular resolution};
\node[label, below=0.5cm of exseq] {Targeted; in-tissue readout};

\end{tikzpicture}
\caption{Overview of spatial transcriptomics technology categories. Sequencing-based
methods (left) capture RNA at spatially barcoded positions and sequence it
off-tissue. Imaging-based methods (center) visualize individual RNA molecules
in situ using fluorescent probes. In situ sequencing methods (right) perform
the sequencing reaction directly within the tissue section.}
\label{fig:spatial:overview}
\end{figure}


% ---------------------------------------------------------------------------
\subsection{Sequencing-Based Spatial Methods}
\label{sec:spatial:seqbased}

These methods capture RNA from intact tissue sections onto spatially barcoded
surfaces. The tissue is permeabilized, RNA molecules are captured by
underlying barcoded oligonucleotides, and the resulting cDNA is sequenced
off-tissue using standard Illumina platforms. The spatial barcode allows each
transcript to be mapped back to its location in the tissue.

\subsubsection{10x Visium}

The most widely adopted spatial transcriptomics platform:

\begin{protocolbox}[title={10x Visium Workflow}]
\begin{enumerate}[nosep]
  \item Mount a fresh-frozen or FFPE tissue section (10\,$\mu$m thick) onto a
    Visium capture slide.
  \item Stain and image the section (H\&E for fresh-frozen; immunofluorescence
    for FFPE).
  \item Permeabilize the tissue to release mRNA (fresh-frozen) or after probe
    hybridization (FFPE).
  \item mRNA binds to spatially barcoded oligo-dT probes printed on the slide
    surface.
  \item Reverse transcription incorporates the spatial barcode into cDNA.
  \item cDNA is collected, amplified, and sequenced on Illumina.
  \item Reads are mapped and demultiplexed by spatial barcode using
    \software{Space Ranger}.
\end{enumerate}
\end{protocolbox}

\textbf{Key specifications:}
\begin{itemize}[nosep]
  \item \textbf{Resolution:} 55\,$\mu$m diameter capture spots arranged in a
    hexagonal grid with 100\,$\mu$m center-to-center spacing. Each spot covers
    1--10 cells (depending on tissue and cell size).
  \item \textbf{Gene coverage:} Genome-wide (no panel required); typically
    $\sim$3,000--8,000 genes detected per spot.
  \item \textbf{Capture area:} 6.5 mm $\times$ 6.5 mm, containing $\sim$5,000
    spots.
  \item \textbf{Sequencing depth:} 50,000--100,000 reads per spot recommended.
  \item \textbf{Sample types:} Fresh-frozen and FFPE (with dedicated Visium
    for FFPE workflow using probe-based capture).
\end{itemize}

\begin{warningbox}
Visium spots are not single cells. Each 55\,$\mu$m spot typically contains
multiple cells, creating a mixture of cell types within each spot. Computational
deconvolution methods (see \cref{sec:spatial:deconv}) are needed to infer
cell type composition per spot.
\end{warningbox}

\subsubsection{10x Visium HD}

The next-generation Visium platform with dramatically improved resolution:

\begin{itemize}[nosep]
  \item \textbf{Resolution:} 2\,$\mu$m $\times$ 2\,$\mu$m barcoded squares,
    providing near-single-cell or subcellular resolution.
  \item \textbf{Binning:} Data can be analyzed at 2\,$\mu$m, 8\,$\mu$m, or
    16\,$\mu$m resolution by binning adjacent squares.
  \item \textbf{Gene coverage:} Genome-wide.
  \item \textbf{Continuous tiling:} No gaps between capture positions, unlike
    the original Visium.
  \item \textbf{Sample types:} FFPE tissues; fresh-frozen support was added
    subsequently.
  \item Analyzed with \software{Space Ranger} (HD mode).
\end{itemize}

\subsubsection{Slide-seq and Slide-seqV2}

\begin{itemize}[nosep]
  \item \textbf{Principle:} A dense monolayer of DNA-barcoded 10\,$\mu$m beads
    (``puck'') is used as the capture surface. Beads are decoded by sequential
    hybridization (\textit{Slide-seq}) or by sequencing (Slide-seqV2) to
    determine their spatial position.
  \item \textbf{Resolution:} 10\,$\mu$m (near single-cell for most tissues).
  \item \textbf{Gene detection:} Lower than Visium ($\sim$500--2,000 genes per
    bead), but higher spatial resolution.
  \item \textbf{Open-source:} Protocols and analysis tools are publicly
    available.
\end{itemize}

\subsubsection{Stereo-seq (BGI)}

\begin{itemize}[nosep]
  \item \textbf{Principle:} DNA nanoball (DNB)-based spatial barcoding on a
    patterned array. Achieves nanoscale resolution ($\sim$500\,nm spot size).
  \item \textbf{Capture area:} Up to full tissue sections (centimeter scale).
  \item \textbf{Resolution:} Subcellular ($\sim$0.5\,$\mu$m), binnable to any
    desired resolution.
  \item \textbf{Notable application:} Used to create the first spatially
    resolved transcriptomic atlas of an entire mouse embryo (MOSTA).
  \item \textbf{Gene coverage:} Genome-wide.
\end{itemize}

\subsubsection{Other Sequencing-Based Methods}

\begin{description}[style=nextline]
  \item[HDST (High-Definition Spatial Transcriptomics)] Uses 2\,$\mu$m
    barcoded wells. An early demonstration of high-resolution capture-based
    spatial transcriptomics.
  \item[Seq-Scope] Repurposes Illumina flow cells as spatial capture surfaces,
    achieving $\sim$0.5--1\,$\mu$m resolution with genome-wide gene coverage
    at very low cost per unit area.
\end{description}


\subsection{Imaging-Based Spatial Methods (In Situ Hybridization)}
\label{sec:spatial:imaging}

Imaging-based methods detect individual RNA molecules in intact tissue by
hybridizing fluorescent probes directly to transcripts. They achieve
subcellular resolution and can image specific panels of 100--10,000+ genes.

\subsubsection{MERFISH (Multiplexed Error-Robust FISH)}

Developed by the Zhuang lab (Harvard) and commercialized by Vizgen:

\begin{keyconcept}[title={How MERFISH Works}]
\begin{enumerate}[nosep]
  \item Design encoding probes for each target gene. Each gene is assigned a
    unique binary barcode (e.g., a 16-bit code using a modified Hamming
    distance error-correcting scheme).
  \item Hybridize all encoding probes simultaneously to the tissue.
  \item Perform multiple rounds of sequential imaging:
    \begin{itemize}[nosep]
      \item In each round, readout probes complementary to specific barcode
        positions bind and fluoresce.
      \item Image the tissue, record fluorescent spots.
      \item Strip the readout probes (photobleach or chemically remove).
      \item Repeat for the next round.
    \end{itemize}
  \item Each RNA molecule generates a pattern of ON/OFF signals across rounds,
    spelling out its binary barcode.
  \item Error correction (Hamming distance codes) allows identification and
    correction of single-bit errors, dramatically reducing false positives.
\end{enumerate}
\end{keyconcept}

\textbf{Key specifications:}
\begin{itemize}[nosep]
  \item \textbf{Gene panel:} 100--1,000+ genes (MERSCOPE platform from Vizgen
    supports panels of 140--1,000 genes; custom panels possible).
  \item \textbf{Resolution:} Single-molecule, subcellular
    ($\sim$100\,nm localization).
  \item \textbf{Throughput:} Millions of cells per tissue section.
  \item \textbf{Imaging time:} 12--48 hours per tissue section depending on
    panel size and area.
\end{itemize}

\subsubsection{seqFISH+}

\begin{itemize}[nosep]
  \item An extension of seqFISH that profiles $>$10,000 genes per cell.
  \item Uses sequential rounds of hybridization with pseudo-color barcoding
    (multiple fluorescent channels per round).
  \item Achieves transcriptome-scale imaging but at lower throughput than
    MERFISH due to the large number of imaging rounds required.
  \item Academic protocol; not commercially available as a turnkey system.
\end{itemize}

\subsubsection{10x Xenium}

The in situ platform from 10x Genomics:

\begin{itemize}[nosep]
  \item \textbf{Principle:} Padlock probe-based detection. Pairs of probes
    hybridize to target RNA; when correctly paired, they are ligated and
    amplified via rolling circle amplification (RCA), creating a bright
    fluorescent spot.
  \item \textbf{Gene panel:} Pre-designed panels (human brain, breast cancer,
    lung, immune, multi-tissue) covering 300--500 genes, plus custom panels
    up to 480 genes. Higher multiplex panels continue to be released.
  \item \textbf{Resolution:} Subcellular (detects individual transcripts within
    cells and can resolve nuclear vs.\ cytoplasmic localization).
  \item \textbf{Throughput:} Large tissue areas (full tissue sections);
    automated instrument with $\sim$24-hour run time per slide.
  \item \textbf{FFPE compatible:} Works with FFPE tissues, enabling retrospective
    spatial analysis of archival samples.
  \item \textbf{Integrated analysis:} Results can be directly overlaid with
    Visium data from adjacent sections for combined genome-wide + targeted
    spatial analysis.
\end{itemize}

\subsubsection{CosMx Spatial Molecular Imager (NanoString)}

\begin{itemize}[nosep]
  \item Probe-based in situ hybridization platform.
  \item Gene panels of 1,000--6,000 genes for RNA and 64--100 proteins
    simultaneously (multi-omic spatial profiling).
  \item Subcellular resolution; automated imaging system.
  \item Compatible with FFPE tissues.
  \item Multi-omic capability (RNA + protein in the same tissue section) is a
    key differentiator.
\end{itemize}

\subsubsection{FISH-Based Principles}

All imaging-based spatial methods rely on variations of fluorescence in situ
hybridization (FISH):

\begin{description}[style=nextline]
  \item[Single-molecule FISH (smFISH)] The foundational method: multiple
    fluorescent probes tile along a single RNA transcript, producing a
    diffraction-limited spot visible under a fluorescence microscope. Detects
    one gene per fluorescence channel (typically 3--5 genes per experiment).
  \item[Combinatorial barcoding] MERFISH and seqFISH assign each gene a
    binary barcode read out across multiple rounds of imaging. This
    exponentially scales the number of genes: $n$ rounds with error correction
    can encode thousands of genes.
  \item[Sequential hybridization] Probes are applied, imaged, removed, and
    replaced in successive rounds. Each round reveals a different subset of
    genes or barcode positions.
\end{description}


\subsection{In Situ Sequencing}
\label{sec:spatial:insitu}

In situ sequencing (ISS) methods perform the sequencing reaction directly
within the tissue, reading out short barcodes that identify each transcript.

\begin{description}[style=nextline]
  \item[ISS (In Situ Sequencing)] Padlock probes hybridize to target RNAs;
    gaps are filled and ligated to form circular DNA. Rolling circle
    amplification (RCA) creates dense DNA nanoballs that are sequenced in situ
    by sequencing-by-ligation. Typically reads 4--5 bases (sufficient for
    barcode identification).
  \item[BaristaSeq] An improved ISS method with higher efficiency gap-filling
    and sequencing-by-synthesis readout. Compatible with tissue clearing for
    3D spatial transcriptomics.
  \item[STARmap / STARmap PLUS] Uses SNAIL (specific amplification of nucleic
    acids via intramolecular ligation) probes for highly multiplexed in situ
    sequencing. STARmap PLUS scales to 1,000+ genes. Compatible with intact
    tissue (hydrogel-embedded) for volumetric 3D spatial transcriptomics.
  \item[ExSeq (Expansion Sequencing)] Combines expansion microscopy (physically
    enlarging the tissue by embedding in a swellable hydrogel) with in situ
    RNA sequencing. Achieves nanoscale resolution of RNA localization within
    subcellular compartments (dendrites, synapses).
\end{description}


% ---------------------------------------------------------------------------
\section{Comparison of Spatial Transcriptomics Platforms}
\label{sec:spatial:comparison}

\begin{table}[htbp]
\centering
\caption{Comprehensive comparison of spatial transcriptomics platforms.}
\label{tab:spatial:platforms}
\footnotesize
\begin{tabularx}{\textwidth}{lXXXXX}
\toprule
\textbf{Feature} & \textbf{Visium} & \textbf{Visium HD} & \textbf{MERFISH} & \textbf{Xenium} & \textbf{Stereo-seq} \\
\midrule
Category & Seq-based & Seq-based & Imaging & Imaging & Seq-based \\
Resolution & 55\,$\mu$m & 2\,$\mu$m & Subcellular & Subcellular & 0.5\,$\mu$m \\
Gene coverage & Genome-wide & Genome-wide & 100--1,000 & 300--500 & Genome-wide \\
Cells per section & $\sim$5,000 spots & $\sim$millions of bins & Millions & Millions & Millions \\
FFPE compatible & Yes (v2) & Yes & Limited & Yes & Yes \\
Instrument & None (slide) & None (slide) & MERSCOPE & Xenium Analyzer & STOmics \\
Run time & 1--2 days (wet lab) & 1--2 days (wet lab) & 12--48 h imaging & $\sim$24 h & 1--2 days \\
Analysis software & Space Ranger & Space Ranger & Vizgen MERSCOPE & Xenium Ranger/Explorer & STOmics tools \\
Cost per section & \$\$\$ & \$\$\$ & \$\$\$\$ & \$\$\$ & \$\$\$ \\
Best use case & Genome-wide discovery & High-res discovery & Targeted, subcellular & Clinical FFPE, targeted & Large-area mapping \\
\bottomrule
\end{tabularx}
\end{table}

\begin{table}[htbp]
\centering
\caption{Extended comparison including additional platforms.}
\label{tab:spatial:platforms_extended}
\footnotesize
\begin{tabularx}{\textwidth}{lXXXXX}
\toprule
\textbf{Feature} & \textbf{Slide-seqV2} & \textbf{seqFISH+} & \textbf{CosMx} & \textbf{STARmap PLUS} & \textbf{Seq-Scope} \\
\midrule
Category & Seq-based & Imaging & Imaging & In situ seq. & Seq-based \\
Resolution & 10\,$\mu$m & Subcellular & Subcellular & Subcellular & $\sim$0.6\,$\mu$m \\
Gene coverage & Genome-wide & 10,000+ & 1,000--6,000 & 1,000+ & Genome-wide \\
Multi-omic & No & No & RNA + protein & No & No \\
Commercial & No (academic) & No (academic) & Yes (NanoString) & No (academic) & No (academic) \\
3D capable & No & Limited & No & Yes & No \\
\bottomrule
\end{tabularx}
\end{table}

\begin{tipbox}
Choosing a spatial technology depends on your biological question:
\begin{itemize}[nosep]
  \item \textbf{Unbiased discovery (genome-wide):} Visium, Visium HD, or
    Stereo-seq.
  \item \textbf{Targeted subcellular resolution:} MERFISH (Vizgen), Xenium
    (10x), or CosMx (NanoString).
  \item \textbf{FFPE clinical samples:} Xenium, Visium FFPE, or CosMx.
  \item \textbf{Combined RNA + protein:} CosMx (NanoString).
  \item \textbf{3D spatial transcriptomics:} STARmap PLUS or ExSeq.
\end{itemize}
Many studies now combine a genome-wide method (e.g., Visium HD) with a targeted
imaging method (e.g., Xenium) on serial sections to get both breadth and
resolution.
\end{tipbox}


% ---------------------------------------------------------------------------
\section{Computational Analysis for Spatial Transcriptomics}
\label{sec:spatial:computation}

Spatial transcriptomics data requires specialized computational methods that
account for the spatial dimensions in addition to the gene expression matrix.

\subsection{Spot Deconvolution}
\label{sec:spatial:deconv}

For Visium and other methods where each measurement captures multiple cells,
deconvolution estimates the proportion of each cell type within each spatial
location by leveraging a single-cell RNA-seq reference:

\begin{description}[style=nextline]
  \item[\software{cell2location}] Bayesian model that estimates absolute cell
    type abundances per spot by decomposing the spatial expression profile into
    a mixture of reference scRNA-seq signatures. Provides uncertainty estimates
    and handles overdispersion. One of the most widely used and best-performing
    methods in benchmarks.
  \item[\software{RCTD} (Robust Cell Type Decomposition)] Part of the
    \software{spacexr} R package. Fits a weighted regression model to estimate
    cell type proportions. Fast and robust; handles platform-specific technical
    differences between spatial and scRNA-seq data.
  \item[\software{Tangram}] Deep learning-based method that maps scRNA-seq
    cells to spatial locations by learning an optimal alignment between the
    two modalities. Can operate at both spot (Visium) and single-cell (Xenium,
    MERFISH) resolution.
  \item[\software{SPOTlight}] Uses non-negative matrix factorization (NMF) to
    decompose Visium spots into cell type proportions.
  \item[\software{stereoscope}] Probabilistic method based on negative binomial
    distributions for deconvolving spatial spots.
\end{description}

\begin{keyconcept}[title={Why Deconvolution Needs a scRNA-seq Reference}]
Deconvolution is fundamentally a source-separation problem: given a mixture
signal (the spatial spot), determine the contributing sources (cell types).
This requires knowing what each source looks like in isolation, which comes
from a scRNA-seq reference of the same tissue. The quality of deconvolution
depends critically on the completeness and accuracy of the reference atlas.
Missing cell types in the reference will be missed in the spatial data.
\end{keyconcept}

\subsubsection{Mathematical Framework for Spot Deconvolution}
\label{subsec:spatial:deconv_math}

The central idea behind methods like \software{RCTD} and
\software{cell2location} is that the observed expression profile at each
spatial location is a mixture of cell-type-specific expression signatures.
Formally, for spot $s$ with observed gene expression vector $Y_s$ (a vector of
length $G$, one entry per gene), the linear mixture model is:
\begin{equation}
  \label{eq:spatial:deconv_linear}
  Y_s = \sum_{k=1}^{K} w_{sk} \cdot \mu_k + \epsilon_s
\end{equation}
where $K$ is the number of cell types in the reference, $\mu_k$ is the
reference expression profile (mean expression vector) for cell type $k$,
$w_{sk} \geq 0$ is the weight (proportion or abundance) of cell type $k$ in
spot $s$, and $\epsilon_s$ represents noise (technical and biological).

\begin{description}[style=nextline]
  \item[\software{RCTD} formulation] Models the observed counts at each spot
    as a weighted sum of cell-type profiles using a platform-corrected
    regression framework. The weights $w_{sk}$ are constrained to be
    non-negative and are estimated via maximum likelihood under a Poisson or
    negative binomial noise model. RCTD explicitly accounts for
    platform-specific effects (differences in capture efficiency and gene
    detection between the spatial and scRNA-seq platforms) by estimating
    platform-specific normalisation factors.
  \item[\software{cell2location} formulation] Uses a hierarchical Bayesian
    model where the expected expression at spot $s$ for gene $g$ is:
    \begin{equation}
      \label{eq:spatial:cell2loc}
      \mathbb{E}[Y_{sg}] = \sum_{k=1}^{K} w_{sk} \cdot g_{kg} \cdot y_s
    \end{equation}
    where $g_{kg}$ is the reference expression of gene $g$ in cell type $k$
    (estimated from the scRNA-seq data), $y_s$ is a spot-level scaling factor,
    and $w_{sk}$ represents the absolute number of cells of type $k$ at spot
    $s$. Bayesian inference (using variational inference in \software{Pyro})
    provides full posterior distributions over all parameters, yielding
    uncertainty estimates for cell type abundances---a key advantage for
    downstream hypothesis testing.
\end{description}

\begin{tipbox}
When choosing a deconvolution method, consider whether you need
\emph{proportions} (what fraction of each spot is cell type $k$?) or
\emph{absolute abundances} (how many cells of type $k$?).
\software{RCTD} estimates proportions by default; \software{cell2location}
estimates absolute abundances. For comparing cell type composition across spots,
proportions suffice. For estimating total cell counts (e.g., immune infiltration
density), absolute abundances are more informative.
\end{tipbox}

\subsection{Spatially Variable Gene Detection}

Identifying genes whose expression varies as a function of spatial location,
beyond what would be expected by chance:

\begin{description}[style=nextline]
  \item[\software{SpatialDE}] Gaussian process regression model that tests
    whether the spatial variance component of gene expression is significant.
    Identifies genes with spatial patterns (e.g., gradients, periodic patterns,
    hotspots) and classifies them by spatial pattern type.
  \item[\software{SPARK}] Spatial Pattern Recognition via Kernels. Tests for
    spatially variable expression using a generalized linear spatial model with
    multiple spatial kernel functions. More powerful than SpatialDE for
    detecting diverse spatial patterns.
  \item[\software{Moran's I}] A classical spatial autocorrelation statistic
    applied to gene expression. Positive Moran's I indicates spatial clustering;
    negative values indicate spatial dispersion. Simple and fast; implemented in
    Squidpy and Giotto.
  \item[\software{nnSVG}] Nearest-neighbor Gaussian process model for
    identifying spatially variable genes. Scales better than SpatialDE to
    high-resolution platforms like Visium HD.
\end{description}

\subsubsection{Moran's I: Spatial Autocorrelation in Detail}
\label{subsec:spatial:morans_i}

Moran's $I$ is the most widely used statistic for quantifying spatial
autocorrelation of gene expression. For a gene with expression values
$x_1, x_2, \ldots, x_n$ measured at $n$ spatial locations, Moran's $I$ is
defined as:
\begin{equation}
  \label{eq:spatial:morans_i}
  I = \frac{n}{\sum_{i} \sum_{j} w_{ij}}
    \cdot
    \frac{\sum_{i} \sum_{j} w_{ij}\,(x_i - \bar{x})(x_j - \bar{x})}
         {\sum_{i} (x_i - \bar{x})^2}
\end{equation}
where $\bar{x} = \frac{1}{n}\sum_i x_i$ is the global mean expression,
$w_{ij}$ is the spatial weight between locations $i$ and $j$ (typically 1 if
$i$ and $j$ are neighbors and 0 otherwise, or a distance-decay function), and
the double sum runs over all pairs of locations.

\textbf{Interpretation:}
\begin{itemize}[nosep]
  \item $I > 0$: \textbf{Positive spatial autocorrelation} (clustering).
    Nearby locations tend to have similar expression values. Values close to
    +1 indicate strong spatial clustering.
  \item $I \approx 0$: \textbf{No spatial autocorrelation} (random spatial
    pattern). Expression is distributed independently of spatial position.
  \item $I < 0$: \textbf{Negative spatial autocorrelation} (dispersion).
    Nearby locations tend to have dissimilar expression values, forming a
    checkerboard-like pattern. This is rare for gene expression data.
\end{itemize}

The expected value under the null hypothesis of spatial randomness is
$\mathbb{E}[I] = -1/(n-1)$, which approaches zero for large $n$.
Statistical significance is assessed either analytically (under normality
assumptions) or by a permutation test that randomly shuffles expression values
across locations and recomputes $I$ to build a null distribution.

\begin{keyconcept}[title={Choosing the Spatial Weight Matrix}]
The spatial weight matrix $W = \{w_{ij}\}$ encodes the definition of ``spatial
proximity'' and critically affects the results:
\begin{itemize}[nosep]
  \item \textbf{Binary adjacency:} $w_{ij} = 1$ if spots $i$ and $j$ are direct
    neighbors (e.g., adjacent hexagons in Visium). Simple and fast.
  \item \textbf{KNN graph:} $w_{ij} = 1$ for the $k$ nearest neighbors of each
    spot. Robust when spot density varies across the tissue.
  \item \textbf{Distance decay:} $w_{ij} = f(d_{ij})$, e.g., Gaussian kernel
    $\exp(-d_{ij}^2 / 2h^2)$ or inverse distance. Captures longer-range spatial
    correlations (e.g., paracrine signaling gradients).
\end{itemize}
In \software{Squidpy}, the default spatial neighbors graph uses Delaunay
triangulation or KNN, and Moran's $I$ is computed via
\texttt{sq.gr.spatial\_autocorr(adata, mode="moran")}.
\end{keyconcept}

\subsection{Spatial Domain Identification}

Identifying tissue regions (``spatial domains'') with coherent gene expression
programs---analogous to clustering in scRNA-seq but incorporating spatial
proximity:

\begin{description}[style=nextline]
  \item[\software{BayesSpace}] Bayesian approach that models spatial
    transcriptomics data on a lattice graph, encouraging neighboring spots to
    belong to the same cluster. Produces spatially smooth domain assignments.
    Supports Visium and Slide-seq data.
  \item[\software{SpaGCN}] Uses a graph convolutional network that integrates
    gene expression, spatial coordinates, and histology image features to
    identify spatial domains. Learns optimal representations that combine
    molecular and morphological information.
  \item[\software{STAGATE}] Graph attention auto-encoder that identifies spatial
    domains by learning a latent representation that incorporates spatial
    neighbor information via attention mechanisms.
  \item[\software{BANKSY}] Builds a neighborhood expression matrix that
    combines each spot's own expression with the average expression of its
    spatial neighbors. Standard clustering on this augmented matrix produces
    spatially coherent domains without specialized spatial algorithms.
\end{description}

\subsection{Integration with scRNA-seq}

Integrating spatial transcriptomics data with scRNA-seq enables cell type
mapping and enhances the resolution of spatial data:

\begin{itemize}[nosep]
  \item \textbf{Deconvolution} (see \cref{sec:spatial:deconv}): Infer cell type
    proportions at each spatial location.
  \item \textbf{Imputation:} Predict the expression of unmeasured genes in
    targeted spatial assays (MERFISH, Xenium) by transferring information from
    a genome-wide scRNA-seq reference. Tools: \software{Tangram},
    \software{gimVI} (scvi-tools), \software{novoSpaRc}.
  \item \textbf{Label transfer:} Project cell type annotations from scRNA-seq
    reference onto spatial spots using anchor-based integration (Seurat) or
    probabilistic assignment (cell2location, Tangram).
\end{itemize}

\subsection{Ligand--Receptor Analysis in Spatial Context}

Spatial data enables direct testing of cell--cell communication hypotheses
because the physical proximity of sender and receiver cells is known:

\begin{description}[style=nextline]
  \item[\software{stLearn}] Combines spatial distance with ligand--receptor
    co-expression to identify spatially constrained cell--cell interactions.
    Only considers interactions between cell types that are physically adjacent.
  \item[\software{COMMOT}] Uses optimal transport theory to model the spatial
    flow of signaling molecules from sender to receiver locations. Produces
    spatial vector fields showing the direction and magnitude of each signaling
    pathway.
  \item[\software{MISTy}] Multi-view framework that models gene expression at
    each spot as a function of its intrinsic state, its immediate neighborhood
    (juxtacrine signaling), and its broader tissue context (paracrine
    signaling). Identifies the spatial scale at which different pathways
    operate.
  \item[\software{Squidpy}] General-purpose spatial analysis framework (see
    below) that includes neighborhood enrichment analysis and ligand--receptor
    co-occurrence testing.
\end{description}

\subsubsection{Ligand--Receptor Communication Scoring: Mathematical Framework}
\label{subsec:spatial:lr_scoring}

Spatial ligand--receptor analysis quantifies the probability or strength of
signaling between spatially proximal cell populations. The \software{CellChat}
framework (originally developed for scRNA-seq) and the spatial extension
\software{COMMOT} both model communication probability based on the expression
levels of ligand and receptor genes, weighted by spatial proximity.

\textbf{CellChat communication probability.} For a ligand--receptor pair
$(L, R)$ between sender cell type $A$ and receiver cell type $B$, CellChat
computes a communication score based on the law of mass action:
\begin{equation}
  \label{eq:spatial:cellchat_prob}
  P_{L,R}^{A \to B} = \frac{\bar{x}_L^A \cdot \bar{x}_R^B}
    {K_h + \bar{x}_L^A \cdot \bar{x}_R^B}
\end{equation}
where $\bar{x}_L^A$ is the average expression of ligand $L$ in sender
population $A$, $\bar{x}_R^B$ is the average expression of receptor $R$ in
receiver population $B$, and $K_h$ is a half-saturation constant (Hill function
parameterization) that prevents extreme scores when expression is very high.
For multi-subunit receptors (e.g., heterodimeric complexes), the receptor
expression is taken as the geometric mean of subunit expression levels.

\textbf{Spatial weighting in COMMOT.} The \software{COMMOT} framework extends
ligand--receptor scoring to spatial transcriptomics by incorporating spatial
distance as a modulating factor. For each pair of spatial locations $(i, j)$
with distance $d_{ij}$, the spatial communication score is:
\begin{equation}
  \label{eq:spatial:commot_score}
  S_{ij}^{L,R} = x_{L,i} \cdot x_{R,j} \cdot K(d_{ij})
\end{equation}
where $x_{L,i}$ is the expression of ligand $L$ at location $i$, $x_{R,j}$ is
the expression of receptor $R$ at location $j$, and $K(d_{ij})$ is a spatial
kernel function that down-weights interactions at greater distances. Common
choices include:
\begin{itemize}[nosep]
  \item \textbf{Cutoff kernel:} $K(d) = \mathbf{1}[d \leq d_{\max}]$ (binary;
    interactions only within a maximum range).
  \item \textbf{Gaussian kernel:} $K(d) = \exp(-d^2 / 2\sigma^2)$ (smooth
    decay with characteristic scale $\sigma$).
\end{itemize}

\software{COMMOT} uses optimal transport theory to aggregate these pairwise
scores into \textbf{spatial communication vector fields}, showing the net
direction and magnitude of signaling flow for each pathway across the tissue.
These vector fields can reveal directional signaling patterns such as morphogen
gradients or immune cell recruitment.

\begin{warningbox}
Ligand--receptor scoring in spatial transcriptomics assumes that mRNA expression
is a reasonable proxy for protein abundance and secretion. This is not always
true: post-transcriptional regulation, protein half-life, and secretion
dynamics can decouple mRNA and protein levels. Spatial proteomics data (e.g.,
from CosMx or CODEX) can provide complementary validation for key signaling
interactions identified by transcriptomic methods.
\end{warningbox}

\subsection{Data Visualization and Analysis Frameworks}

\begin{description}[style=nextline]
  \item[\software{Squidpy}] Python package for spatial omics analysis built on
    top of Scanpy/AnnData. Provides image processing (segmentation, feature
    extraction), graph construction (spatial neighbors), spatial statistics
    (Moran's I, co-occurrence, centrality scores), and interactive
    visualization. The primary Python framework for spatial transcriptomics.
  \item[\software{Giotto}] R and Python package providing a comprehensive suite
    for spatial data analysis: preprocessing, spatial network construction,
    spatial gene detection, domain identification, cell--cell interaction, and
    visualization with interactive viewers.
  \item[\software{Seurat} (spatial module)] Seurat v5 includes native support
    for spatial data (Visium, Xenium, MERFISH, Slide-seq). Functions for
    spatial feature plots, integration with scRNA-seq, and spatial
    dimensionality reduction.
  \item[\software{Voyager}] R/Bioconductor package implementing spatial
    statistics methods (variograms, correlograms, Moran's I, Geary's C) for
    spatial transcriptomics, built on the SpatialFeatureExperiment data
    structure.
  \item[\software{napari} + \software{spatialdata}] Python imaging viewer
    (napari) combined with the spatialdata framework for multi-modal spatial
    data handling (images, coordinates, annotations, expression). Enables
    interactive exploration of spatial datasets.
\end{description}


% ---------------------------------------------------------------------------
\section{Tissue Handling and Sample Preparation}
\label{sec:spatial:tissue}

Spatial transcriptomics requires careful tissue preparation---often more
demanding than standard sequencing experiments because tissue morphology must
be preserved.

\subsection{Tissue Sectioning}

\begin{itemize}[nosep]
  \item \textbf{Fresh-frozen (FF):} Tissue is flash-frozen in OCT compound
    (optimal cutting temperature) or isopentane. Sectioned at 10\,$\mu$m on a
    cryostat. Best RNA quality but limited long-term storage.
  \item \textbf{FFPE:} Formalin-fixed, paraffin-embedded. Sectioned at
    5--10\,$\mu$m on a microtome. RNA is crosslinked and fragmented but
    preserved indefinitely. Compatible with Visium FFPE, Xenium, CosMx.
  \item \textbf{Section quality:} Sections must be uniform in thickness,
    wrinkle-free, and fully adherent to the capture area. Tissue folds,
    tears, or bubbles create artifacts in spatial data.
\end{itemize}

\subsection{Permeabilization Optimization}

For sequencing-based methods (Visium), the tissue must be permeabilized to
release mRNA while keeping it localized above its capture spot.

\begin{protocolbox}[title={Visium Permeabilization Optimization}]
A tissue optimization (TO) experiment is performed before the main experiment:
\begin{enumerate}[nosep]
  \item Mount serial sections on a Visium TO slide.
  \item Permeabilize each section for a different duration (e.g., 3, 6, 12, 18,
    24, 30 minutes).
  \item Perform reverse transcription with fluorescent nucleotides.
  \item Image the fluorescent cDNA footprint under a microscope.
  \item Select the permeabilization time that maximizes signal (cDNA
    fluorescence) while maintaining spatial resolution (sharp edges, no
    lateral diffusion).
\end{enumerate}
Typical optimal permeabilization time: 6--18 minutes for most tissues, but
varies widely (e.g., brain requires less, skeletal muscle requires more).
\end{protocolbox}

\begin{warningbox}
Over-permeabilization causes RNA to diffuse laterally, blurring spatial
resolution. Under-permeabilization results in poor RNA capture. Each tissue
type requires its own optimization. Do not skip the tissue optimization step.
\end{warningbox}

\subsection{Imaging}

\begin{itemize}[nosep]
  \item \textbf{H\&E staining:} Standard histological stain for tissue
    morphology. Used with fresh-frozen Visium. The H\&E image is registered to
    the spatial gene expression data, allowing correlation of transcriptomic
    patterns with morphological features.
  \item \textbf{Immunofluorescence (IF):} Antibody-based staining for specific
    proteins. Used with FFPE Visium and Xenium. Provides protein-level
    information that complements the transcriptomic data.
  \item \textbf{Cell segmentation:} For imaging-based methods, individual cells
    must be segmented from the tissue image. Methods include DAPI-based nuclear
    segmentation (watershed, Cellpose, StarDist) and membrane staining
    (phalloidin, WGA). Accurate segmentation is critical for assigning
    transcripts to individual cells.
\end{itemize}


% ---------------------------------------------------------------------------
\section{Emerging Directions: Spatial Multi-Omics}
\label{sec:spatial:emerging}

\subsection{Spatial Epigenomics}

\begin{description}[style=nextline]
  \item[Spatial ATAC-seq] Methods such as spatial-ATAC-seq and
    Spatial-CUT\&Tag map chromatin accessibility and histone modifications
    in tissue sections. These reveal spatially resolved regulatory landscapes
    (enhancer usage, TF binding) that complement spatial gene expression.
  \item[10x Visium + ATAC] Emerging protocols from 10x Genomics that combine
    gene expression and chromatin accessibility on the same Visium slide.
\end{description}

\subsection{Spatial Proteomics}

\begin{description}[style=nextline]
  \item[CODEX (CO-Detection by indEXing)] Iterative antibody staining and
    imaging that profiles 40--60 proteins at single-cell resolution in intact
    tissues. Uses DNA-conjugated antibodies and fluorescent reporters that
    are sequentially activated and quenched.
  \item[MIBI-TOF (Multiplexed Ion Beam Imaging)] Uses metal-conjugated
    antibodies and time-of-flight mass spectrometry for highly multiplexed
    protein imaging ($>$40 markers).
  \item[CosMx multi-omic:] Simultaneous RNA and protein detection on the same
    tissue section (see above).
\end{description}

\subsection{Spatial Multi-Omic Integration}

The frontier of spatial biology is measuring multiple molecular layers
(transcriptome, epigenome, proteome, metabolome) in the same tissue section
or in serial sections from the same block:

\begin{itemize}[nosep]
  \item Serial section approaches: alternate sections for different assays
    (e.g., Visium + Xenium + CODEX on serial 10\,$\mu$m sections).
  \item Same-section multi-omics: CosMx (RNA + protein), Visium + IF, and
    emerging combined ATAC + RNA spatial methods.
  \item Computational integration: Tools such as \software{MEFISTO} (multi-omic
    factor analysis in space and time) and \software{MOFA+} enable joint
    analysis of spatial multi-omic datasets.
\end{itemize}


% ---------------------------------------------------------------------------
\section{Snakemake Workflow for Spatial Transcriptomics}
\label{sec:spatial:snakemake}

\begin{lstlisting}[language=Python, caption={Snakemake workflow for 10x Visium
spatial transcriptomics analysis.}, label={lst:spatial:snakemake}]
# Snakefile for Spatial Transcriptomics (10x Visium)
# ====================================================
# Pipeline: Space Ranger -> Squidpy / Scanpy analysis

import pandas as pd

configfile: "config/config.yaml"

samples = pd.read_csv(config["samples"], sep="\t")
SAMPLES = samples["sample"].tolist()

REFERENCE = config["spaceranger_ref"]

rule all:
    input:
        expand("results/spaceranger/{sample}/outs/"
               "filtered_feature_bc_matrix.h5",
               sample=SAMPLES),
        "results/squidpy/spatial_analysis.h5ad",
        "results/squidpy/spatial_domains.pdf",
        "results/deconvolution/cell2location_results.h5ad"

rule spaceranger_count:
    input:
        fastqs="data/raw/{sample}/",
        image="data/images/{sample}.tif",
        reference=REFERENCE
    output:
        h5="results/spaceranger/{sample}/outs/"
           "filtered_feature_bc_matrix.h5",
        spatial="results/spaceranger/{sample}/outs/spatial/",
        web="results/spaceranger/{sample}/outs/"
            "web_summary.html"
    threads: 12
    resources:
        mem_mb=64000
    params:
        outdir="results/spaceranger/{sample}",
        slide=lambda wc: samples.loc[
            samples["sample"] == wc.sample,
            "slide_serial"].iloc[0],
        area=lambda wc: samples.loc[
            samples["sample"] == wc.sample,
            "capture_area"].iloc[0]
    shell:
        """
        spaceranger count \
          --id={wildcards.sample} \
          --transcriptome={input.reference} \
          --fastqs={input.fastqs} \
          --image={input.image} \
          --slide={params.slide} \
          --area={params.area} \
          --localcores={threads} \
          --localmem=60 \
          --output-dir={params.outdir}
        """

rule squidpy_analysis:
    input:
        h5_files=expand(
            "results/spaceranger/{sample}/outs/"
            "filtered_feature_bc_matrix.h5",
            sample=SAMPLES),
        spatial_dirs=expand(
            "results/spaceranger/{sample}/outs/spatial/",
            sample=SAMPLES)
    output:
        h5ad="results/squidpy/spatial_analysis.h5ad",
        domains="results/squidpy/spatial_domains.pdf",
        svg_genes="results/squidpy/spatially_variable_genes.csv",
        co_occur="results/squidpy/co_occurrence.pdf"
    conda: "envs/squidpy.yaml"
    script:
        "scripts/run_squidpy.py"

rule cell2location:
    input:
        spatial="results/squidpy/spatial_analysis.h5ad",
        sc_ref=config["scrna_reference"]
    output:
        results="results/deconvolution/"
                "cell2location_results.h5ad",
        plots="results/deconvolution/"
              "cell_type_maps.pdf"
    resources:
        mem_mb=32000
    conda: "envs/cell2location.yaml"
    script:
        "scripts/run_cell2location.py"

rule spatial_communication:
    input:
        h5ad="results/squidpy/spatial_analysis.h5ad"
    output:
        commot="results/communication/"
               "spatial_signaling.h5ad",
        plots="results/communication/"
              "signaling_vectors.pdf"
    conda: "envs/commot.yaml"
    script:
        "scripts/run_commot.py"
\end{lstlisting}


% ---------------------------------------------------------------------------
\section{Nextflow Workflow for Spatial Transcriptomics}
\label{sec:spatial:nextflow}

\begin{lstlisting}[language=Java, caption={Nextflow DSL2 workflow for 10x
Visium spatial transcriptomics analysis.},
label={lst:spatial:nextflow}]
#!/usr/bin/env nextflow
// Spatial Transcriptomics Pipeline (Nextflow DSL2)
// ==================================================
// Pipeline: Space Ranger -> Squidpy -> cell2location

nextflow.enable.dsl = 2

params.sample_sheet = "config/samples.tsv"
params.reference    = "resources/refdata-gex-GRCh38-2024-A"
params.scrna_ref    = "resources/scrna_reference.h5ad"
params.outdir       = "results"

log.info """
==========================================
 Spatial Transcriptomics Pipeline
==========================================
 sample_sheet : ${params.sample_sheet}
 reference    : ${params.reference}
 scRNA ref    : ${params.scrna_ref}
 outdir       : ${params.outdir}
==========================================
"""

// Parse sample sheet
Channel
    .fromPath(params.sample_sheet)
    .splitCsv(header: true, sep: '\t')
    .map { row ->
        tuple(
            row.sample,
            file(row.fastq_dir),
            file(row.image),
            row.slide_serial,
            row.capture_area
        )
    }
    .set { samples_ch }

process SPACERANGER_COUNT {
    tag "$sample_id"
    publishDir "${params.outdir}/spaceranger/${sample_id}",
        mode: 'copy'
    cpus 12
    memory '64 GB'

    input:
    tuple val(sample_id), path(fastqs), path(image),
          val(slide), val(area)
    path(reference)

    output:
    tuple val(sample_id),
        path("outs/filtered_feature_bc_matrix.h5"),
        path("outs/spatial/"),
        emit: counts
    path("outs/web_summary.html"), emit: web_summary

    script:
    """
    spaceranger count \
      --id=${sample_id} \
      --transcriptome=${reference} \
      --fastqs=${fastqs} \
      --image=${image} \
      --slide=${slide} \
      --area=${area} \
      --localcores=${task.cpus} \
      --localmem=60
    """
}

process SQUIDPY_ANALYSIS {
    publishDir "${params.outdir}/squidpy", mode: 'copy'
    cpus 4
    memory '32 GB'

    input:
    path(h5_files)
    path(spatial_dirs)

    output:
    path("spatial_analysis.h5ad"),          emit: h5ad
    path("spatial_domains.pdf"),            emit: domains
    path("spatially_variable_genes.csv"),   emit: svg
    path("co_occurrence.pdf"),              emit: co_occur

    script:
    """
    python ${projectDir}/bin/run_squidpy.py \
      --input_h5 ${h5_files} \
      --spatial_dirs ${spatial_dirs} \
      --outdir .
    """
}

process CELL2LOCATION {
    publishDir "${params.outdir}/deconvolution", mode: 'copy'
    memory '32 GB'
    accelerator 1

    input:
    path(spatial_h5ad)
    path(scrna_ref)

    output:
    path("cell2location_results.h5ad"), emit: results
    path("cell_type_maps.pdf"),         emit: plots

    script:
    """
    python ${projectDir}/bin/run_cell2location.py \
      --spatial ${spatial_h5ad} \
      --sc_reference ${scrna_ref} \
      --outdir .
    """
}

process SPATIAL_COMMUNICATION {
    publishDir "${params.outdir}/communication", mode: 'copy'
    memory '16 GB'

    input:
    path(spatial_h5ad)

    output:
    path("spatial_signaling.h5ad"),  emit: results
    path("signaling_vectors.pdf"),   emit: plots

    script:
    """
    python ${projectDir}/bin/run_commot.py \
      --input ${spatial_h5ad} \
      --outdir .
    """
}

workflow {
    reference = Channel.fromPath(params.reference)

    // Space Ranger quantification
    SPACERANGER_COUNT(samples_ch, reference.collect())

    // Collect all samples for joint analysis
    SQUIDPY_ANALYSIS(
        SPACERANGER_COUNT.out.counts
            .map { it[1] }.collect(),
        SPACERANGER_COUNT.out.counts
            .map { it[2] }.collect()
    )

    // Cell type deconvolution
    CELL2LOCATION(
        SQUIDPY_ANALYSIS.out.h5ad,
        Channel.fromPath(params.scrna_ref)
    )

    // Spatial communication
    SPATIAL_COMMUNICATION(
        SQUIDPY_ANALYSIS.out.h5ad
    )
}
\end{lstlisting}

\begin{tipbox}
For production spatial transcriptomics analysis, the \software{nf-core/spatialtranscriptomics}
pipeline provides a Nextflow workflow supporting Space Ranger, quality control,
normalization, clustering, and basic spatial analysis with standardized outputs.
\end{tipbox}


% ---------------------------------------------------------------------------
\section{Practical Experimental Design Considerations}
\label{sec:spatial:design}

\subsection{Choosing the Right Technology}

\begin{itemize}[nosep]
  \item \textbf{Discovery vs.\ hypothesis-driven:} If you need genome-wide
    unbiased discovery, use a sequencing-based method (Visium, Visium HD,
    Stereo-seq). If you have a defined gene panel, imaging-based methods
    (MERFISH, Xenium) provide higher resolution and single-cell quantification.
  \item \textbf{Sample type:} FFPE samples have the widest compatibility
    (Visium FFPE, Xenium, CosMx). Fresh-frozen tissue generally yields better
    RNA quality for sequencing-based methods.
  \item \textbf{Resolution requirements:} If single-cell resolution is essential,
    use imaging-based methods or Visium HD. Standard Visium (55\,$\mu$m)
    requires computational deconvolution for cell type analysis.
  \item \textbf{Throughput:} For large cohorts (tens to hundreds of samples),
    consider the cost per section and instrument availability. Visium slides
    are relatively straightforward; imaging instruments have limited throughput
    (one section per $\sim$24 hours).
\end{itemize}

\subsection{Replication and Sample Size}

\begin{itemize}[nosep]
  \item Biological replicates (n $\geq$ 3 per condition) are needed for
    statistical testing of spatial gene expression differences.
  \item Technical replicates (serial sections from the same block) assess
    within-sample variability.
  \item Consider tissue heterogeneity: one section may not represent the entire
    organ. Multiple sections from different regions may be needed.
\end{itemize}

\subsection{Combining Spatial and Single-Cell Data}

\begin{protocolbox}[title={Recommended Multi-Modal Experimental Design}]
For comprehensive tissue profiling, a combined approach is often optimal:
\begin{enumerate}[nosep]
  \item \textbf{scRNA-seq:} Generate a reference atlas of cell types from
    dissociated tissue (10x Chromium, 5,000--10,000 cells per sample, n $\geq$ 3).
  \item \textbf{Spatial transcriptomics (genome-wide):} Profile serial sections
    with Visium or Visium HD for unbiased spatial gene expression mapping.
  \item \textbf{Spatial transcriptomics (targeted):} Profile adjacent sections
    with Xenium or MERFISH for single-cell resolution of key pathways.
  \item \textbf{Integration:} Use cell2location or Tangram to project scRNA-seq
    cell types onto spatial data; impute unmeasured genes in targeted data.
\end{enumerate}
This multi-modal approach provides genome-wide coverage (Visium), single-cell
resolution (Xenium), and a comprehensive cell type reference (scRNA-seq).
\end{protocolbox}


% ---------------------------------------------------------------------------
\section{Chapter Summary}
\label{sec:spatial:summary}

Spatial transcriptomics represents one of the most exciting frontiers in
genomics, bridging the gap between molecular profiling and tissue biology. In
this chapter we covered:

\begin{itemize}[nosep]
  \item Why spatial context matters: tissue architecture, cell neighborhoods,
    gradients, and the limitations of dissociation-based methods.
  \item Sequencing-based spatial methods: 10x Visium (55\,$\mu$m spots), Visium HD
    (2\,$\mu$m), Slide-seqV2 (10\,$\mu$m beads), Stereo-seq (nanoscale), and
    Seq-Scope.
  \item Imaging-based spatial methods: MERFISH (Vizgen), seqFISH+, 10x Xenium,
    and CosMx (NanoString), including the principles of combinatorial barcoding
    and sequential hybridization.
  \item In situ sequencing methods: ISS, BaristaSeq, STARmap PLUS, and ExSeq.
  \item Comprehensive platform comparisons covering resolution, gene coverage,
    throughput, and cost.
  \item Computational analysis: spot deconvolution (cell2location, RCTD,
    Tangram), spatially variable gene detection (SpatialDE, SPARK, Moran's I),
    spatial domain identification (BayesSpace, SpaGCN, STAGATE), scRNA-seq
    integration, and ligand--receptor analysis in spatial context (stLearn,
    COMMOT, MISTy).
  \item Visualization and analysis frameworks: Squidpy, Giotto, Seurat spatial
    module, Voyager.
  \item Tissue handling: sectioning (FF and FFPE), permeabilization
    optimization, imaging, and cell segmentation.
  \item Emerging directions: spatial ATAC-seq, spatial proteomics (CODEX,
    MIBI-TOF), and multi-omic spatial integration.
  \item Practical workflows in both Snakemake and Nextflow for Visium analysis.
\end{itemize}

Spatial transcriptomics, combined with the single-cell methods described in
\cref{ch:scrnaseq} and the bulk approaches in \cref{ch:rnaseq}, provides a
complete toolkit for transcriptomic analysis at every scale---from
population-level averages to individual molecules in their native tissue
context.
